% Generated draft: requires mathematical and source review.
\documentclass[11pt]{article}
\usepackage[margin=25mm]{geometry}
\usepackage{fontspec}
\setmainfont{Latin Modern Roman}
\usepackage{amsmath,amssymb,mathtools,unicode-math}
\setmathfont{Latin Modern Math}
\usepackage{enumitem,needspace,xurl,hyperref}
\setlength{\parindent}{0pt}
\setlength{\parskip}{5pt}
\setlength{\emergencystretch}{3em}
\newcommand{\N}{\mathbb{N}}
\newcommand{\Z}{\mathbb{Z}}
\newcommand{\Q}{\mathbb{Q}}
\newcommand{\R}{\mathbb{R}}
\newcommand{\C}{\mathbb{C}}
\newcommand{\F}{\mathbb{F}}
\newcommand{\A}{\mathbb{A}}
\newcommand{\PP}{\mathbb{P}}
\newcommand{\E}{\mathbb{E}}
\newcommand{\eps}{\varepsilon}
\newcommand{\dd}{\,\mathrm{d}}
\begin{document}
\section*{UnsolvedMath ID: 30003834}
\textbf{Source problem code:} catalogue entry
\subsection{Definitions and mathematical statement}
Let \(\mathbb H^2\) denote the hyperbolic plane, and let \(\Delta\) be a tessellation of \(\mathbb H^2\). The supplied record imposes no further conditions on the tessellation. In the usual countable-set convention, the finitary symmetric group is
\[
S_\infty=\{\sigma\in\operatorname{Sym}(\Omega):\{x\in\Omega:\sigma(x)\ne x\}\text{ is finite}\},
\]
where \(\Omega\) is a countably infinite set. The record does not specify the underlying set or its relation to \(\Delta\).

Thompson's group \(V\) is the group of homeomorphisms of \(\{0,1\}^{\N}\) given by finite binary prefix replacements: partition the space into finitely many cylinder sets \(u_i\{0,1\}^{\N}\), choose another such partition \(v_i\{0,1\}^{\N}\), and map \(u_iw\) to \(v_iw\).

Write \(G=\phi(\Delta)\). The construction \(\phi(\Delta)\) is not defined in the supplied record. Its quotient notation presupposes a specified normal subgroup of \(G\) identified with \(S_\infty\). Subject to that missing construction and identification, the hypothesis is
\[
G/S_\infty\cong V.
\]
\textbf{Question.} For every tessellation covered by the intended construction and satisfying this hypothesis, do \(G\) and \(G/S_\infty\) have the same finiteness properties?

The expression ``same finiteness properties'' is not specified in the record. For example, a group \(H\) is of type \(F_n\), for an integer \(n\geq0\), if it admits a classifying CW complex \(K(H,1)\) with finite \(n\)-skeleton. If the intended properties are precisely these geometric finiteness properties, the question becomes whether
\[
\forall n\in\Z_{\geq0},\qquad
G\text{ is of type }F_n
\ \Longleftrightarrow\
G/S_\infty\text{ is of type }F_n.
\]
This is a conditional interpretation, not a resolution of the source's terminology. Homological finiteness properties, such as type \(FP_n\) with specified coefficients, may also be intended. The missing definition of \(\phi(\Delta)\), the admissible tessellations, the normal subgroup identification, and the intended family of finiteness properties require clarification before a fully determined statement is available.
\subsection{Short English statement}
The supplied question asks whether a group associated with a hyperbolic tessellation has the same finiteness properties as its quotient by finitary permutations, assuming that quotient is Thompson's group \(V\). The construction of the group and the intended meaning of finiteness properties remain unspecified.
\subsection{Sources}
\begin{itemize}
\item Supplied local UnsolvedMath registry record, permanent ID \(30003834\), source problem code \texttt{catalogue entry}, titled ``Finiteness Properties of Hyperbolic Tessellation Groups,'' in the collection ``Oberwolfach Reports Open Problems.'' No individual proposer is supplied; the proposed year is recorded as 2018.
\item Supplied source citation: \emph{Cohomological and Metric Properties of Groups of Homeomorphisms of R} (2019), Oberwolfach Reports, pp.\ 1579--1633, DOI \texttt{10.4171/owr/2018/26}. \url{https://doi.org/10.4171/owr/2018/26}. The citation's 2019 label and the 2018 year in the DOI are retained as supplied.
\item Supplied literature source labeled ``Oberwolfach Report 2019'': \url{https://ems.press/content/serial-article-files/46748?nt=1}.
\item Supplied literature source attributed to Farley--Hughes: \url{https://doi.org/10.4171/JCA/91}. No article title or pinpoint supporting this question is supplied.
\item Literature assessment dated \texttt{2026-08-21}: the supplied curation labels the question open, describes equality of finiteness properties as open in general, and reports known special Euclidean and Houghton-type cases. This is supplied dated curation derived from the OWR record, not an independently verified literature assessment. The record explicitly requests independent verification of time-sensitive status and cited claims.
\end{itemize}
\end{document}
